\documentclass[11pt]{article}
\usepackage[T1]{fontenc}
\usepackage[textwidth=250pt]{geometry}
\usepackage{amsmath,amssymb}
\usepackage{microtype}
\pagestyle{empty}
\begin{document}
Suppose that $a+b=c$ and that $u=a+1$ and $v=a+2b$ hold for all integers.
Observe that $u$ and $v$ are coprime whenever $\gcd(a,b)=1$, so that the
pair $(u,v)$ satisfies $u\le v$ and $v-u=2b-1$ in every case we check.

If $x<y$ then $x+z<y+z$ for every real $z$, and the relation $x\ne y$ holds
whenever $x<y$ does. Moreover $2x<2y$, so that the inequality survives the
scaling by any positive factor $k>0$ that we choose to apply to both sides.

Set $u=a+1$ and $v=a+2b$ for part (c). Observe that $u-v=1-2b$ is odd and
that $2v=2a+4b$ is even, so that $u$ and $v$ never coincide, and the claim
follows at once from these two small observations about parity.
\end{document}
